\begin{lemma}
\label{lemma-relative-colimit-general}
Let $I$ be a directed set. Let $g_i : X_i \to S_i$ be an inverse system of
morphisms of schemes over $I$. Assume $g_i$ is quasi-compact and
quasi-separated and for $i' \geq i$ the transition morphisms
$f_{i'i} : X_{i'} \to X_i$ and $h_{i'i} : S_{i'} \to S_i$ are affine.
Let $g : X \to S$ be the limit of the morphisms $g_i$, see
Limits, Section \ref{limits-section-limits}.
Denote $f_i : X \to X_i$ and $h_i : S \to S_i$ the projections.
Let $(\mathcal{F}_i, \varphi_{i'i})$ be a system of sheaves
on $(X_i, f_{i'i})$. Set $\mathcal{F} = \colim f_i^{-1}\mathcal{F}_i$. Then
$$
R^p g_* \mathcal{F} =
\colim_{i \in I} h_i^{-1}R^p g_{i, *} \mathcal{F}_i
$$
for all $p \geq 0$.
\end{lemma}

\begin{proof}
How is the map of the lemma constructed?
For $i' \geq i$ we have a commutative diagram
$$
\xymatrix{
X \ar[r]_{f_{i'}} \ar[d]_g &
X_{i'} \ar[r]_{f_{i'i}} \ar[d]_{g_{i'}} &
X_i \ar[d]^{g_i} \\
S \ar[r]^{h_{i'}} &
S_{i'} \ar[r]^{h_{i'i}} &
S_i
}
$$
If we combine the base change map
$h_{i'i}^{-1}Rg_{i, *}\mathcal{F}_i \to Rg_{i', *}f_{i'i}^{-1}\mathcal{F}_i$
(Cohomology on Sites, Lemma
\ref{sites-cohomology-lemma-base-change-map-flat-case} or
Remark \ref{sites-cohomology-remark-base-change})
with the map $Rg_{i', *}\varphi_{i'i}$, then we obtain
$\psi_{i'i} : h_{i' i}^{-1} R^p g_{i, *} \mathcal{F}_i \to
R^pg_{i', *} \mathcal{F}_{i'}$. Similarly, using the left square
in the diagram we obtain maps
$\psi_i : h_i^{-1}R^pg_{i, *}\mathcal{F}_i \to R^pg_*\mathcal{F}$.
The maps $h_{i'}^{-1}\psi_{i'i}$ and $\psi_i$ are the maps used in
the statement of the lemma. For this to make sense, we have to check that
$\psi_{i''i} = \psi_{i''i'} \circ h_{i''i'}^{-1}\psi_{i'i}$ and
$\psi_{i'} \circ h_{i'}^{-1}\psi_{i'i} = \psi_i$; this follows
from Cohomology on Sites, Remark
\ref{sites-cohomology-remark-compose-base-change-horizontal}.

\medskip\noindent
Proof of the equality. First proof using
dimension shifting\footnote{You can also use this method
to produce the maps in the lemma.}. For any
$U$ affine and \'etale over $X$ by Theorem \ref{theorem-colimit}
we have
$$
g_*\mathcal{F}(U) =
H^0(U \times_S X, \mathcal{F}) =
\colim H^0(U_i \times_{S_i} X_i, \mathcal{F}_i) =
\colim g_{i, *}\mathcal{F}_i(U_i)
$$
where the colimit is over $i$ large enough such that
there exists an $i$ and $U_i$ affine \'etale over $S_i$
whose base change is $U$ over $S$ (see
Lemma \ref{lemma-colimit-affine-sites}).
The right hand side is equal to
$(\colim h_i^{-1}g_{i, *}\mathcal{F}_i)(U)$ by
Sites, Lemma \ref{sites-lemma-colimit}.
This proves the lemma for $p = 0$.
If $(\mathcal{G}_i, \varphi_{i'i})$ is a system with
$\mathcal{G} = \colim f_i^{-1}\mathcal{G}_i$
such that $\mathcal{G}_i$ is an injective abelian sheaf on $X_i$
for all $i$, then for any $U$ affine and \'etale over $X$ by
Theorem \ref{theorem-colimit} we have
$$
H^p(U \times_S X, \mathcal{G}) =
\colim H^p(U_i \times_{S_i} X_i, \mathcal{G}_i) = 0
$$
for $p > 0$ (same colimit as before). Hence $R^pg_*\mathcal{G} = 0$
and we get the result for $p > 0$ for such a system.
In general we may choose a short exact sequence of systems
$$
0 \to (\mathcal{F}_i, \varphi_{i'i}) \to
(\mathcal{G}_i, \varphi_{i'i}) \to
(\mathcal{Q}_i, \varphi_{i'i}) \to 0
$$
where $(\mathcal{G}_i, \varphi_{i'i})$ is as above, see
Cohomology on Sites, Lemma
\ref{sites-cohomology-lemma-colim-sites-injective}.
By induction the lemma holds for $p - 1$ and by the above we have
vanishing for $p$ and $(\mathcal{G}_i, \varphi_{i'i})$.
Hence the result for $p$
and $(\mathcal{F}_i, \varphi_{i'i})$ by the long exact sequence
of cohomology.

\medskip\noindent
Second proof.
Recall that $S_{affine, \etale} = \colim (S_i)_{affine, \etale}$, see
Lemma \ref{lemma-colimit-affine-sites}. Thus if $U$ is an object of
$S_{affine, \etale}$, then we can write $U = U_i \times_{S_i} S$
for some $i$ and some $U_i$ in $(S_i)_{affine, \etale}$ and
$$
(\colim_{i \in I} h_i^{-1}R^p g_{i, *} \mathcal{F}_i)(U) =
\colim_{i' \geq i} (R^p g_{i', *}\mathcal{F}_{i'})(U_i \times_{S_i} S_{i'})
$$
by Sites, Lemma \ref{sites-lemma-colimit} and the construction of the
transition maps in the system described above. Since
$R^pg_{i', *}\mathcal{F}_{i'}$ is the sheaf associated to the presheaf
$U_{i'} \mapsto H^p(U_{i'} \times_{S_{i'}} X_{i'}, \mathcal{F}_{i'})$
and since $R^pg_*\mathcal{F}$ is the sheaf associated to the presheaf
$U \mapsto H^p(U \times_S X, \mathcal{F})$
(Lemma \ref{lemma-higher-direct-images})
we obtain a canonical commutative diagram
$$
\xymatrix{
\colim_{i' \geq i}
H^p(U_i \times_{S_i} X_{i'}, \mathcal{F}_{i'}) \ar[r] \ar[d] &
\colim_{i' \geq i}
(R^p g_{i', *}\mathcal{F}_{i'})(U_i \times_{S_i} S_{i'}) \ar[d] \\
H^p(U \times_S X, \mathcal{F}) \ar[r] &
R^pg_*\mathcal{F}(U)
}
$$
Observe that the left hand vertical arrow is an isomorphism
by Theorem \ref{theorem-colimit}. We're trying to show that
the right hand vertical arrow is an isomorphism. However, we
already know that the source and target of this arrow
are sheaves on $S_{affine, \etale}$. Hence it suffices to
show: (1) an element in the target, locally comes from an
element in the source and (2) an element in the source
which maps to zero in the target locally vanishes.
Part (1) follows immediately from the above and the fact that
the lower horizontal arrow comes from a map of presheaves
which becomes an isomorphism after sheafification.
For part (2), say $\xi \in  \colim_{i' \geq i}
(R^p g_{i', *}\mathcal{F}_{i'})(U_i \times_{S_i} S_{i'})$
is in the kernel. Choose an $i' \geq i$ and
$\xi_{i'} \in (R^p g_{i', *}\mathcal{F}_{i'})(U_i \times_{S_i} S_{i'})$
representing $\xi$.
Choose a standard \'etale covering
$\{U_{i', k} \to U_i \times_{S_i} S_{i'}\}_{k = 1, \ldots, m}$
such that $\xi_{i'}|_{U_{i', k}}$ comes from
$\xi_{i', k} \in H^p(U_{i', k} \times_{S_{i'}} X_{i'}, \mathcal{F}_{i'})$.
Since it is enough to prove that $\xi$ dies locally, we
may replace $U$ by the members of the \'etale
covering $\{U_{i', k} \times_{S_{i'}} S \to U = U_i \times_{S_i} S\}$.
After this replacement we see that $\xi$ is the image of
an element $\xi'$ of the group
$\colim_{i' \geq i} H^p(U_i \times_{S_i} X_{i'}, \mathcal{F}_{i'})$
in the diagram above. Since $\xi'$ maps to zero in $R^pg_*\mathcal{F}(U)$
we can do another replacement and assume that $\xi'$ maps
to zero in $H^p(U \times_S X, \mathcal{F})$.
However, since the left vertical arrow is an isomorphism
we then conclude $\xi' = 0$ hence $\xi = 0$ as desired.
\end{proof}